\chapter{Mappe concettuali}
\label{cap:mappe}

% ---------------------------------------------------------------------
\section{Quale norma si applica?}

Prima delle due competenze richieste dal corso (lezione del 02/11). Le domande \textbf{non sono
alternative}: vanno poste tutte, perché le norme si cumulano.

\begin{center}
\begin{tikzpicture}[
  node distance=6mm and 8mm,
  q/.style={draw, rounded corners, fill=blue!6, text width=6.6cm, align=left,
    font=\small, inner sep=5pt, minimum height=1.7cm},
  r/.style={draw, thick, fill=green!8, text width=5.4cm, align=left,
    font=\small, inner sep=5pt, minimum height=1.7cm},
  >=Stealth]
\node[q] (q1) {Sono coinvolti \textbf{dati personali} di persone fisiche?};
\node[r, right=of q1] (r1) {\textbf{GDPR}\\ se data breach: artt.~33--34};

\node[q, below=of q1] (q2) {Il soggetto è un'\textbf{entità finanziaria} (banca, istituto di
  pagamento, assicurazione, \ldots) o un suo fornitore ICT?};
\node[r, right=of q2] (r2) {\textbf{DORA}\\ lex specialis rispetto a NIS2; ai fornitori ICT si
  applica per via contrattuale};

\node[q, below=of q2] (q3) {Il soggetto è \textbf{essenziale o importante} in un settore degli
  Allegati I--II di NIS2?};
\node[r, right=of q3] (r3) {\textbf{NIS2}\\ in Italia D.Lgs.~138/2024, autorità ACN};

\node[q, below=of q3] (q4) {È un \textbf{prodotto con elementi digitali} (hardware o software
  installabile, connesso) messo a disposizione sul mercato UE e non coperto da un regime verticale?};
\node[r, right=of q4] (r4) {\textbf{CRA}\\ obblighi di segnalazione dall'11/9/2026; tutto dall'11/12/2027};

\node[q, below=of q4] (q5) {È o contiene un \textbf{sistema di IA}?};
\node[r, right=of q5] (r5) {\textbf{AI Act}\\ alto rischio: presunzione di conformità all'art.~15 tramite CRA};

\node[q, below=of q5] (q6) {Si vuole dimostrare la sicurezza con una \textbf{certificazione UE}?};
\node[r, right=of q6] (r6) {\textbf{CSA}\\ schemi europei di certificazione, ENISA};

\foreach \i in {1,...,6} {\draw[->] (q\i) -- node[above, font=\scriptsize]{sì} (r\i);}
\foreach \i/\j in {1/2,2/3,3/4,4/5,5/6} {\draw[->] (q\i) -- node[left, font=\scriptsize]{in ogni caso} (q\j);}
\end{tikzpicture}
\end{center}

\begin{nota}
\begin{itemize}
  \item Ospedale, società energetica, telco: NIS2 e, se ci sono dati personali, GDPR.
  \item Banca o istituto di pagamento: DORA e GDPR (NIS2 cede a DORA dove le norme si sovrappongono).
  \item Fabbricante di dispositivi IoT o di app: CRA e, quasi sempre, GDPR.
  \item Un certificato ISO 27001 o CSA \textbf{supporta} l'accountability ma \textbf{non prova} la
    conformità legale.
\end{itemize}
\end{nota}

% ---------------------------------------------------------------------
\section{Il percorso in caso di incidente}

Seconda competenza richiesta dal corso: individuare il percorso corretto quando si verifica un incidente.

\begin{center}
\begin{tikzpicture}[
  node distance=4mm,
  s/.style={draw, rounded corners, fill=orange!8, text width=12.5cm, align=left,
    font=\small, inner sep=5pt},
  >=Stealth]
\node[s] (a) {\textbf{1. Rilevazione ed escalation}: chi riceve il segnale (SOC, reperibile, call
  center) lo porta subito a chi decide, compresi DPO e legale.};
\node[s, below=of a] (b) {\textbf{2. Contenimento e preservazione delle evidenze}: isolamento, immagini
  forensi, log, catena di custodia.};
\node[s, below=of b] (c) {\textbf{3. Qualificazione}: quali norme sono coinvolte (mappa precedente)?
  Data breach? Incidente significativo/grave? Vulnerabilità attivamente sfruttata?};
\node[s, below=of c] (d) {\textbf{4. Consapevolezza}: fissare e \textbf{documentare} il momento di
  avvio di ciascun orologio, che può essere diverso per ciascuna norma.};
\node[s, below=of d] (e) {\textbf{5. Valutazione della severità}: ENISA per il GDPR; criteri di
  significatività/gravità per NIS2, DORA e CRA.};
\node[s, below=of e] (f) {\textbf{6. Notifiche in parallelo} verso ciascuna autorità, rispettando
  ogni scadenza (Capitolo~\ref{cap:scadenze}).};
\node[s, below=of f] (g) {\textbf{7. Comunicazioni}: interessati (art.~34 GDPR), utenti (CRA),
  destinatari dei servizi e pubblico (NIS2), retailer, stampa. Trasparenza: niente smentite false.};
\node[s, below=of g] (h) {\textbf{8. Registri e chiusura}: registro degli incidenti e registro dei
  data breach allineati; report finali; lezioni apprese e aggiornamento di risk assessment e DPIA.};
\foreach \x/\y in {a/b,b/c,c/d,d/e,e/f,f/g,g/h} {\draw[->] (\x) -- (\y);}
\end{tikzpicture}
\end{center}

\dacompletare{Altre mappe da aggiungere man mano: ruoli GDPR (titolare, responsabile,
sub-responsabile); classificazione dei prodotti CRA; classificazione per rischio dell'AI Act;
soggetti NIS2.}
